% !TeX root = ../../tesis.tex
%%%%%%%%%%%%%%%%%%%%%%%%%%%%%%%%%%%%%%%%%%%%%%%%%%%%%%%%%%%%%%%%%%%
% SECCIÓN 5.4.2 — PENALIZACIÓN POR TRANSFERENCIA (Tb)
%%%%%%%%%%%%%%%%%%%%%%%%%%%%%%%%%%%%%%%%%%%%%%%%%%%%%%%%%%%%%%%%%%%
\subsection{Penalización por Transferencia (\texorpdfstring{$T_b$}{Tb})}
\label{subsec:fase2-tb}

Cambiar de línea o de sistema no es gratuito en términos de tiempo.
Más allá de la caminata entre andenes, el usuario debe esperar la
llegada del siguiente servicio: una espera que depende de la frecuencia
del sistema al que transborda. La penalización por transferencia ($T_b$)
incorpora ese costo al grafo asignando un peso adicional a cada arista
de tipo trasbordo.

\begin{equation}
    T_b = \frac{F}{2}
    \label{eq:penalizacion-transferencia}
\end{equation}

\noindent donde $F$ es la frecuencia del sistema destino expresada en
minutos entre servicios. El factor $\tfrac{1}{2}$ representa la
espera media esperada bajo llegadas aleatorias al andén: en promedio,
el usuario llega en el punto medio del intervalo entre dos servicios
consecutivos.

La penalización actúa exclusivamente sobre las 20,482 aristas de
trasbordo peatonal del grafo —aquellas que conectan nodos de sistemas
distintos dentro de un mismo complejo de transferencia— y no modifica
el tiempo de las aristas de servicio real. Con $T_b$ aplicado, el
grafo refleja no solo el tiempo de traslado sobre cada línea, sino el
costo acumulado de cada cambio de sistema que la ruta óptima imponga
al usuario.

Ambos parámetros —$CF$ y $T_b$— quedan incorporados al grafo ponderado
antes de ejecutar el algoritmo de ruta mínima. La Fase~3 opera sobre
ese grafo para derivar el tiempo de viaje promedio de la red y su
distribución territorial.
